\documentclass[11pt]{article}
% ===========================================================================
%  Shared preamble for the AP Calculus AB exam-preparation worksheets.
%  Input right after \documentclass. Files beginning with "_" are not compiled
%  on their own. Navy and gold, to match the course's lesson pages.
% ===========================================================================
\usepackage[letterpaper,top=2.3cm,bottom=2.2cm,left=1.7cm,right=1.7cm,headheight=28pt,footskip=30pt]{geometry}
\usepackage{amsmath,amssymb}
\usepackage[T1]{fontenc}
\usepackage[utf8]{inputenc}
\usepackage{lmodern}
\usepackage{xcolor}
\usepackage[most]{tcolorbox}
\usepackage{enumitem}
\usepackage{fancyhdr}
\usepackage{array}
\usepackage{tikz}
\usepackage{pgfplots}
\pgfplotsset{compat=1.18}

\definecolor{navy}{HTML}{161E54}
\definecolor{ink}{HTML}{27357E}
\definecolor{gold}{HTML}{B98A00}
\definecolor{goldbg}{HTML}{FFFDF3}
\definecolor{inkbg}{HTML}{F6F7FD}
\definecolor{rule}{HTML}{C9CFF0}
\definecolor{okgreen}{HTML}{1B7A44}
\definecolor{okbg}{HTML}{F1F8F3}

\pagestyle{fancy}
\fancyhf{}
\fancyhead[L]{\footnotesize Name: \underline{\hspace{4cm}}}
\fancyhead[C]{\textbf{\color{navy}Integration Academy}\\[-2pt]{\scriptsize Calculus AB \textperiodcentered{} AP\textsuperscript{\textregistered} Exam Preparation}}
\fancyfoot[L]{\tiny\color{black!55}\parbox[t]{0.78\textwidth}{AP\textsuperscript{\textregistered} is a registered trademark of the College Board, which was not involved in the production of, and does not endorse, this material. \textcopyright{} Integration Academy.}}
\fancyfoot[R]{\scriptsize\thepage}
\renewcommand{\headrulewidth}{1.2pt}
\renewcommand{\headrule}{\hbox to\headwidth{\color{gold}\leaders\hrule height \headrulewidth\hfill}}
\renewcommand{\footrulewidth}{0pt}
\setlength{\parindent}{0pt}
\setlength{\parskip}{4pt}

% Title block: \wstitle{code}{title}{unit line}
\newcommand{\wstitle}[3]{\begin{center}{\small\bfseries\color{gold}WORKSHEET #1}\\[3pt]{\LARGE\bfseries\color{navy}#2}\\[3pt]{\small\color{black!60}#3}\end{center}\vspace{2pt}}
% Section heading
\newcommand{\sech}[1]{\par\vspace{10pt}{\large\bfseries\color{navy}#1}\par\nobreak\vspace{2pt}{\color{rule}\hrule height0.8pt}\vspace{7pt}}
% A level tag printed beside a group of questions
\newcommand{\level}[1]{\par\vspace{6pt}{\footnotesize\bfseries\color{gold}\MakeUppercase{#1}}\par\nobreak\vspace{2pt}}

\newtcolorbox{ideabox}{colframe=ink,colback=inkbg,coltitle=white,colbacktitle=ink,fonttitle=\bfseries\small,title={Key ideas},arc=2pt,boxrule=0.8pt,left=8pt,right=8pt}
\newtcolorbox{example}[1]{colframe=rule,colback=white,coltitle=navy,colbacktitle=inkbg,fonttitle=\bfseries\small,title={#1},arc=2pt,boxrule=0.8pt,borderline west={3pt}{0pt}{ink},left=10pt,right=8pt}
\newtcolorbox{frbox}[1]{colframe=gold,colback=goldbg,coltitle=white,colbacktitle=gold,fonttitle=\bfseries\small,title={#1},arc=2pt,boxrule=0.8pt,breakable,left=8pt,right=8pt}
\newtcolorbox{solbox}{colframe=okgreen,colback=okbg,coltitle=white,colbacktitle=okgreen,fonttitle=\bfseries\small,title={Worked solutions},arc=2pt,boxrule=0.8pt,breakable,left=8pt,right=8pt}
\newtcolorbox{scorebox}[1]{colframe=navy,colback=white,coltitle=white,colbacktitle=navy,fonttitle=\bfseries\small,title={#1},arc=2pt,boxrule=0.8pt,breakable,left=8pt,right=8pt}

\newcommand{\soln}{\par\smallskip\textbf{\color{ink}Solution.}\ }
\newcommand{\result}[1]{\par\smallskip\textbf{#1}}
% blank working space of a given height
\newcommand{\work}[1]{\par\vspace{#1}}
% a point value in the margin of a scoring guide
\newcommand{\pts}[1]{\hfill{\footnotesize\bfseries\color{gold}[#1]}}

% Coordinate grid for graph questions: \begin{apgraph}{xmin}{xmax}{ymin}{ymax} ... \end{apgraph}
\newenvironment{apgraph}[4]{\begin{center}\begin{tikzpicture}\begin{axis}[
  axis lines=middle, xmin=#1, xmax=#2, ymin=#3, ymax=#4,
  grid=both, grid style={black!12}, xtick distance=1, ytick distance=1,
  ticklabel style={font=\scriptsize}, xlabel={$x$}, ylabel={$y$},
  width=10cm, height=7.2cm, samples=80, clip=false,
  every axis plot/.append style={line width=1.1pt, color=ink}]}{\end{axis}\end{tikzpicture}\end{center}}
\newcommand{\closeddot}[2]{\addplot[only marks, mark=*, mark size=2.2pt, color=ink] coordinates {(#1,#2)};}
\newcommand{\opendot}[2]{\addplot[only marks, mark=*, mark size=2.2pt, color=ink, mark options={fill=white}] coordinates {(#1,#2)};}

% Multiple-choice options: four across, or stacked when the options are long.
\newcommand{\opts}[4]{\par\vspace{3pt}\begin{tabular}{@{}p{0.235\linewidth}p{0.235\linewidth}p{0.235\linewidth}p{0.235\linewidth}@{}}(A)\ #1 & (B)\ #2 & (C)\ #3 & (D)\ #4\end{tabular}\par\vspace{9pt}}
\newcommand{\optsv}[4]{\par\vspace{3pt}\begin{tabular}{@{}l@{\ \ }p{0.9\linewidth}@{}}(A) & #1\\(B) & #2\\(C) & #3\\(D) & #4\end{tabular}\par\vspace{9pt}}

\fancyhead[R]{\footnotesize\color{black!60}Worksheet 5.1}
\begin{document}

\wstitle{5.1}{The Mean Value Theorem}{Unit 5: Analysing Functions}

\begin{ideabox}
\begin{itemize}[leftmargin=*,itemsep=2pt,topsep=2pt]
  \item \textbf{Mean Value Theorem.} If $f$ is continuous on $[a,b]$ and differentiable on $(a,b)$, there is a $c$ in $(a,b)$ with $f'(c)=\dfrac{f(b)-f(a)}{b-a}$.
  \item In words: at some point the instantaneous rate equals the average rate.
  \item If also $f(a)=f(b)$, then $f'(c)=0$ for some $c$ (Rolle's Theorem).
  \item To \emph{use} the theorem in a justification, state both conditions. \textquotedblleft{}Differentiable, and therefore continuous\textquotedblright{} covers both.
\end{itemize}
\end{ideabox}

\sech{Worked examples}
\begin{example}{Example 1 \textperiodcentered{} Finding c}
Find the value of $c$ guaranteed by the Mean Value Theorem for $f(x)=x^2-2x$ on $[1,4]$.\soln $\dfrac{f(4)-f(1)}{3}=\dfrac{8-(-1)}{3}=3$. Solve $2c-2=3$.\result{$c=\dfrac52$.}
\end{example}

\begin{example}{Example 2 \textperiodcentered{} A reciprocal}
Find the value of $c$ guaranteed by the theorem for $f(x)=\dfrac1x$ on $[1,9]$.\soln $\dfrac{\frac19-1}{8}=-\dfrac19$. Solve $-\dfrac{1}{c^2}=-\dfrac19$ with $1<c<9$.\result{$c=3$.}
\end{example}

\begin{example}{Example 3 \textperiodcentered{} From values}
A differentiable function $g$ has $g(0)=5$ and $g(4)=13$. What does the theorem guarantee?\soln $\dfrac{13-5}{4}=2$, and $g$ is differentiable, hence continuous, on $[0,4]$.\result{There is a $c$ in $(0,4)$ with $g'(c)=2$.}
\end{example}

\newpage
\sech{Practice}

\level{Core}
\begin{enumerate}[leftmargin=*,itemsep=2pt]
  \item Find the value of $c$ guaranteed by the Mean Value Theorem for $f(x)=x^2+1$ on $[0,2]$.\work{2.3cm}
  \item Find the value of $c$ guaranteed by the theorem for $f(x)=x^3$ on $[-1,2]$.\work{2.3cm}
  \item State the two conditions of the Mean Value Theorem.\work{2.3cm}
  \item Explain why the theorem does not apply to $f(x)=|x|$ on $[-1,1]$.\work{2.3cm}
\end{enumerate}

\level{Exam level}
\begin{enumerate}[leftmargin=*,itemsep=2pt,start=5]
  \item Find the value of $c$ guaranteed by the theorem for $f(x)=\sqrt{x}$ on $[1,9]$.\work{2.8cm}
  \item Find the value of $c$ guaranteed by the theorem for $f(x)=\ln x$ on $[1,e]$.\work{2.8cm}
  \item A car travels 150 km along a road in exactly 2 hours. Explain why its speedometer must have read 75 km/h at some instant. State any assumption you make.\work{2.8cm}
  \item A differentiable function $h$ has $h(2)=3$ and $h(7)=3$. Must there be a $c$ in $(2,7)$ with $h'(c)=0$? Must there be a $c$ with $h(c)=0$?\work{2.8cm}
\end{enumerate}

\begin{frbox}{Question 9 \textperiodcentered{} Free-response practice}
Let $f(x)=x^3-x^2$.
\begin{enumerate}[label=(\alph*),leftmargin=*,itemsep=2pt]
  \item Explain why the Mean Value Theorem applies to $f$ on $[0,2]$.\work{2.6cm}
  \item Find the value of $c$ in $(0,2)$ guaranteed by the theorem.\work{2.6cm}
  \item Show that $f(0)=f(1)$, and find the value of $c$ in $(0,1)$ at which $f'(c)=0$.\work{2.6cm}
\end{enumerate}
\end{frbox}

\newpage
\begin{solbox}
\begin{enumerate}[leftmargin=*,itemsep=4pt]
  \item Average rate $\dfrac{5-1}{2}=2$; $2c=2$ gives $c=1$.
  \item Average rate $\dfrac{8-(-1)}{3}=3$; $3c^2=3$ gives $c=1$ (the value $c=-1$ is not inside the interval).
  \item $f$ is continuous on $[a,b]$ and differentiable on $(a,b)$.
  \item $f$ is not differentiable at $x=0$, which lies in $(-1,1)$.
  \item Average rate $\dfrac{3-1}{8}=\dfrac14$; $\dfrac{1}{2\sqrt c}=\dfrac14$ gives $c=4$.
  \item Average rate $\dfrac{1-0}{e-1}$; $\dfrac1c=\dfrac{1}{e-1}$ gives $c=e-1$.
  \item Assuming the position is a differentiable function of time, its average rate of change is $\dfrac{150}{2}=75$ km/h, so by the Mean Value Theorem the velocity equals 75 km/h at some instant.
  \item Yes to the first: $h$ is differentiable and $h(2)=h(7)$, so the theorem gives $h'(c)=\dfrac{3-3}{5}=0$. No to the second: nothing forces $h$ to take the value 0 (for example $h$ could be constant at 3).
  \item \begin{enumerate}[label=(\alph*),leftmargin=*,itemsep=3pt]
    \item $f$ is a polynomial, so it is continuous on $[0,2]$ and differentiable on $(0,2)$.
    \item $\dfrac{f(2)-f(0)}{2}=\dfrac{4}{2}=2$. $3c^2-2c=2$ gives $c=\dfrac{1+\sqrt7}{3}\approx1.215$ (the other root is negative).
    \item $f(0)=0$ and $f(1)=1-1=0$. $f'(c)=3c^2-2c=c(3c-2)=0$ in $(0,1)$ gives $c=\dfrac23$.
  \end{enumerate}
\end{enumerate}
\end{solbox}

\end{document}
